% -------------------------------------------------------------------
% 2. Типографика.
% -------------------------------------------------------------------
\lecture[2026]{2}{Основы типографики}

\begin{attentionbox}{Что знать}
	Гарнитура, шрифт, начертание, кегль, классификация шрифтов
	(антиква, гротеск, акциденция), длина строки, интерлиньяж,
	трекинг, кернинг, «висячая» пунктуация, спецсимволы и~типографская
	раскладка, редактура
\end{attentionbox}

Текст~"--- такой же элемент дизайна, как и~изображения. Пунктуация
по правилам русского языка повышает доверие к~вашему тексту и~к~вам
как к~профессионалу. Чтобы расставлять правильные тире и~кавычки
(«ёлочки» и~„лапки“), пользуйтесь
\href{https://ilyabirman.ru/typography-layout/}{раскладкой Бирмана}
и~\href{%
	https://www.figma.com/community/plugin/907888805660596065/sbol-typograph
}{плагином~SBOL.} 


Расстояния лучше подбирать по правилу внутреннего и~внешнего
(о~нём~"--- ниже), а~для сочетания шрифтов нужно знать их типологии
и~подбирать референсы. Не~растягивайте строку больше 6--8 слов.
Длина строки взаимозависит от~гарнитуры, кегля и~интерлиньяжа,
поэтому универсальные рекомендации вроде 120\% работают не~всегда.
Это означает, что вышеупомянутая рекомендация в~6--8 слов тоже
не~универсальна и~не~подойдёт для узких колонок или текста,
набранного рукописным шрифтом (пожалуйста, не~набирайте основной
текст рукописным шрифтом!)

\begin{definition}
	\tdef{Гарнитура}~"--- семейство шрифтов (Roboto), \tdef{шрифт}~"--- конкретное
	начертание (Roboto Bold Italic). Другие дизайнеры и~так вас
	поймут, но~разницу знать ст\'{о}ит.
\end{definition}

\textbf{Правило внутреннего и~внешнего:} расстояния внутри объекта
должны быть меньше расстояний снаружи. Это частный случай общей
теории близости~"--- ещё одного ключевого правила в~дизайне.
Напоминаем: объекты, расположенные близко друг к~другу,
воспринимаются связанно.

\begin{example}
	Меняя одно из~расстояний, пересматривайте всю иерархию в~текстовом
	блоке:
	\[
		\begin{array}{c}
			\text{внутрибуквенное} \\
			\big\Downarrow \\
			\text{межбуквенное} \\
			\big\Downarrow \\
			\text{пробел} \\
			\big\Downarrow \\
			\text{межстрочное} \\
			\big\Downarrow \\
			\text{между абзацами} \\
			\big\Downarrow \\
			\text{абзацы и~заголовок}
		\end{array}
	\]
\end{example}

Внутрибукенные и~межбуквенные расстояния уже заданы в~шрифте.
В~общем случае дизайнеры полагают, что шрифтовикам виднее, поэтому
менять межбуквенные расстояния основного текста не рекомендуется.
А~вот крупным акцидентным заголовкам часто нужен трекинг~"--- равное
изменение межбуквенного расстояния между всеми буквами в~строке. Это
нужно, потому что у~больших заголовков межбуквенные пробелы
становятся более заметными и~занимают много места. Если попался
некачественный шрифт, пригодится и~кернинг~"--- точечное изменение
расстояния между двумя буквами. Поиграться с~кернингом можно
\href{https://type.method.ac/}{в~игре Kern~Type}. В~живых проектах
заголовок советуем отзеркалить или повернуть вверх ногами. Так проще
увидеть, где нужно компенсировать расстояния.

Чтобы изменить кегль и~посмотреть, как устроен элемент на~«живой»
веб"=странице, пользуйтесь инспектором элементов. Вызывайте окно
клавишей~\keys{F12} или из~контекстного меню «Исследовать элемент»
по правой кнопке мыши. В~панели браузера вы увидите код страницы
и~свойства каждого элемента. 

\begin{codebox}[%
		Базовый HTML~"--- чтобы разбираться, как устроены
		веб-страницы с~помощью инспектора элементов%
	]
	\inputminted{html}{code/html_1.html}
\end{codebox}

Чтобы увеличить или уменьшить кегль, найдите нужный абзац текста
(правая кнопка мыши~"--- «Исследовать элемент») и~измените
CSS-свойство~\cssinline{font-size}. Ваши изменения локальны, поэтому
при~обновлении страницы значение сбросится. Если вы знаете HTML
и~CSS на~базовом уровне, вам проще общаться с~разработчиками
и~контролировать вёрстку ваших макетов. 

Здесь~же советуем ознакомиться с~термином
\href{https://bureau.ru/about/fff/}{«флексить»}~"---
отказываться от~функций, сохраняя полезное действие продукта
при~фиксированном времени и~бюджете. Флекс
и~\href{%
	https://www.artlebedev.ru/kovodstvo/sections/167/
}{прогрессивный джипег}
помогут довести до~финала любой проект. 


% vim:ts=2:sw=2:tw=68
